\documentclass[11pt]{article}
\usepackage[margin=1in]{geometry}
\usepackage{amsmath,amssymb,amsthm}
\usepackage{enumitem}
\usepackage[hidelinks]{hyperref}

\newtheorem{theorem}{Theorem}
\newtheorem{proposition}[theorem]{Proposition}
\newtheorem{lemma}[theorem]{Lemma}
\newtheorem{corollary}[theorem]{Corollary}
\theoremstyle{definition}
\newtheorem{definition}[theorem]{Definition}
\theoremstyle{remark}
\newtheorem{remark}[theorem]{Remark}

\newcommand{\F}{\mathbb{F}}
\newcommand{\Z}{\mathbb{Z}}
\newcommand{\GL}{\mathrm{GL}}

\title{Route-chart lifts over finite rings:\\
$N(4n,5)\ge 144\,n^5$ for every $n\not\equiv 2 \pmod 4$}
\author{Ryosuke Mizuno}
\date{November 2026}

\begin{document}
\maketitle

\begin{abstract}
The route-chart lift of arXiv:2606.15860 combines a $4$-regular
\emph{controller} on $144$ vertices with a \emph{universal route chart} on
$\F_q^5$ and gives graphs of maximum degree $4q$ and diameter at most $5$
on $144\,q^5$ vertices. The paper gives universal charts for $q=3$ and $q=4$.
We show that the construction works over any finite commutative ring that
carries a universal chart, and that such rings exist in every order
$n\not\equiv 2\pmod 4$. Hence $N(4n,5)\ge 144\,n^5$ for every $n\ge 1$ with
$n\not\equiv 2 \pmod 4$, in particular for every prime power $q\ge 3$. We also
show that universal charts exist over every finite field $\F_q$ with $q\ge 3$,
and that none exists over $\F_2$. The first statement is formalized in Lean~4
with Mathlib, including all finite checks, which the Lean kernel verifies.
\end{abstract}

\section{Results}

Let $N(\Delta,D)$ be the largest order of a graph with maximum degree at most
$\Delta$ and diameter at most $D$.

\begin{theorem}\label{thm:main}
For every integer $n\ge 1$ with $n\not\equiv 2\pmod 4$ there is a graph with
$144\,n^5$ vertices and maximum degree at most $4n$ in which every two
vertices are joined by a walk of length at most $5$. Hence
\[
  N(4n,5)\ \ge\ 144\,n^5 \qquad (n\ge 1,\ n\not\equiv 2 \pmod 4).
\]
\end{theorem}

Every prime power $q\ge 3$ is odd or divisible by $4$, so:

\begin{corollary}\label{cor:pp}
$N(4q,5)\ge 144\,q^5$ for every prime power $q\ge 3$.
\end{corollary}

For $n\ge 2$ the order $144\,n^5$ exceeds the Moore bound
$\Delta^4-2\Delta^3+2\Delta^2+1<\Delta^4=256\,n^4$ for maximum degree
$\Delta=4n$ and diameter $4$, so these graphs have diameter exactly $5$.
For $n=3$ and $n=4$ they are the graphs of arXiv:2606.15860 on $34{,}992$ and
$147{,}456$ vertices. For $n=5$ the graph has degree $20$ and $450{,}000$ vertices.
Its edge list was checked directly, by breadth-first search from every vertex
(Section~\ref{sec:verify}).

The construction in arXiv:2606.15860 is stated for finite fields. Our second
result says that the field version also works for every $q\ge 3$.

\begin{theorem}\label{thm:field}
For every prime power $q\ge 3$ there is a universal route chart on $\F_q^5$
(Definition~\ref{def:chart}). With the $144$-vertex controller, the route-chart lift over
$\F_q$ is a $4q$-regular graph of order $144\,q^5$ and diameter $5$.
\end{theorem}

\begin{proposition}\label{prop:two}
There is no universal route chart on $\F_2^5$. Consequently no finite
commutative ring of order $n\equiv 2\pmod 4$ carries a universal route chart,
and Theorem~\ref{thm:main} is the full range of orders reached by route-chart
lifts over finite commutative rings.
\end{proposition}

Theorem~\ref{thm:main} is proved in Sections~\ref{sec:lift}--\ref{sec:proof} and is
formalized in Lean (Section~\ref{sec:lean}). Theorem~\ref{thm:field} and
Proposition~\ref{prop:two} are proved in Sections~\ref{sec:field}
and~\ref{sec:two}. They rely on further computer checks, not on the Lean
formalization.

\section{Route charts over commutative rings}\label{sec:lift}

Throughout, $R$ is a finite commutative ring. Let $S=\{a,a^{-1},b,b^{-1}\}$,
with the involution $t\mapsto t^{-1}$. A word $t_1t_2t_3t_4t_5$ over $S$ is
\emph{reduced} if $t_{i+1}\ne t_i^{-1}$ for $1\le i\le 4$; there are
$4\cdot 3^4=324$ reduced words of length $5$.

\begin{definition}[Controller]
A \emph{controller} on a finite set $C$ is a family of maps
$\sigma_t\colon C\to C$ ($t\in S$) with
$\sigma_{t^{-1}}(\sigma_t(u))=u$ for all $t$ and $u$. It is
\emph{exact-NB-$5$} if for all $u,v\in C$ there is a reduced word
$t_1\cdots t_5$ with $\sigma_{t_5}\cdots\sigma_{t_1}(u)=v$.
\end{definition}

The controller of arXiv:2606.15860, \S A.5, is a simple $4$-regular graph on
$144$ vertices. Its arcs are labelled by $S$ so that every vertex has
exactly one outgoing arc with each label and the reverse of an arc with label
$t$ has label $t^{-1}$ (\S A.3). Writing $\sigma_t(u)$ for the head of the arc
with label $t$ leaving $u$ gives a controller in the above sense. A walk is
nonbacktracking exactly when its label word is reduced. The paper verifies
that this controller is exact-NB-$5$, and we verify it again
(Section~\ref{sec:verify}).

\begin{definition}[Route chart]\label{def:chart}
A \emph{route chart} on $R^5$ assigns to every $t\in S$ a matrix
$A_t\in M_5(R)$ and a vector $b_t\in R^5$ with
\begin{equation}\label{eq:A1}
  A_{t^{-1}}A_t = I, \qquad b_{t^{-1}} = A_{t^{-1}}\,b_t \qquad (t\in S).
\end{equation}
For a word $w=t_1\cdots t_5$ let
\begin{equation}\label{eq:A2}
  M_w=\bigl[\,A_{t_5}A_{t_4}A_{t_3}A_{t_2}b_{t_1}\ \big|\ A_{t_5}A_{t_4}A_{t_3}b_{t_2}\ \big|\
  A_{t_5}A_{t_4}b_{t_3}\ \big|\ A_{t_5}b_{t_4}\ \big|\ b_{t_5}\,\bigr]\in M_5(R).
\end{equation}
The chart is \emph{universal} if $M_w$ has a right inverse for every reduced
word $w$ of length $5$. Over a commutative ring this means that $\det M_w$
is a unit, and over a field that $M_w$ is nonsingular, as in
arXiv:2606.15860.
\end{definition}

\begin{definition}[Lift graph]
Given a controller on $C$ and a route chart on $R^5$, the graph $G$ has vertex
set $C\times R^5$. Two distinct vertices $(u,x)$ and $(v,y)$ are adjacent if
$v=\sigma_t(u)$ and $y=A_tx+\lambda b_t$ for some $t\in S$ and $\lambda\in R$.
\end{definition}

This relation is symmetric. If $v=\sigma_t(u)$ and $y=A_tx+\lambda b_t$, then
$\sigma_{t^{-1}}(v)=u$ and, by~\eqref{eq:A1},
$A_{t^{-1}}y+(-\lambda)b_{t^{-1}}=x+\lambda A_{t^{-1}}b_t-\lambda b_{t^{-1}}=x$.
When the controller is a simple graph and $R$ is a field, this is the
route-chart lift graph of arXiv:2606.15860, \S A.4.

\begin{lemma}\label{lem:lift}
\begin{enumerate}[label=(\alph*)]
\item Every vertex of $G$ has at most $4|R|$ neighbours.
\item If the controller is exact-NB-$5$ and the chart is universal, every two
vertices of $G$ are joined by a walk of length at most $5$.
\end{enumerate}
\end{lemma}

\begin{proof}
(a) The neighbours of $(u,x)$ lie in
$\{(\sigma_t(u),A_tx+\lambda b_t): t\in S,\ \lambda\in R\}$.

(b) Let $(u,x)$ and $(v,y)$ be vertices. Choose a reduced word
$w=t_1\cdots t_5$ with $\sigma_{t_5}\cdots\sigma_{t_1}(u)=v$, and a right
inverse $N$ of $M_w$. Put $A_w=A_{t_5}\cdots A_{t_1}$ and
$\lambda=N(y-A_wx)\in R^5$. Define $x_0=x$, $u_0=u$ and, for
$i=1,\dots,5$, $u_i=\sigma_{t_i}(u_{i-1})$ and
$x_i=A_{t_i}x_{i-1}+\lambda_ib_{t_i}$. Expanding,
\[
  x_5=A_wx+\sum_{j=1}^5\lambda_j\,A_{t_5}\cdots A_{t_{j+1}}b_{t_j}
     =A_wx+M_w\lambda=A_wx+(y-A_wx)=y,
\]
and $u_5=v$. Consecutive vertices $(u_{i-1},x_{i-1})$ and $(u_i,x_i)$ are
equal or adjacent, so deleting repetitions from
$(u_0,x_0),\dots,(u_5,x_5)$ gives a walk of length at most $5$ from $(u,x)$
to $(v,y)$.
\end{proof}

\section{Transport, products and the orders reached}

\begin{lemma}\label{lem:transport}
Let $f\colon R\to R'$ be a homomorphism of commutative rings. Applying $f$ to
all entries maps a route chart on $R^5$ to a route chart on $R'^5$, maps each
$M_w$ to the corresponding matrix of the image chart, and preserves
universality.
\end{lemma}

\begin{proof}
Entrywise application of $f$ commutes with sums and products of matrices and
vectors, so~\eqref{eq:A1}, \eqref{eq:A2} and the equations $M_wN=I$ are preserved.
\end{proof}

\begin{lemma}\label{lem:prod}
If $R_1$ and $R_2$ carry universal route charts, so does $R_1\times R_2$.
\end{lemma}

\begin{proof}
Pair the two charts entrywise: $A_t=(A_t^{(1)},A_t^{(2)})$ and
$b_t=(b_t^{(1)},b_t^{(2)})$. All identities hold componentwise, and the pair
$(N^{(1)},N^{(2)})$ of right inverses is a right inverse of $M_w$.
\end{proof}

Call $n$ \emph{good} if some commutative ring with $n$ elements carries a
universal route chart. The zero ring shows that $1$ is good, since every
matrix over it is the identity. By Lemma~\ref{lem:prod}, if $m$ and $n$ are
good then so is $mn$.

\begin{lemma}\label{lem:reduce}
If $4$, $8$ and every odd prime are good, then every $n\ge 1$ with
$n\not\equiv 2\pmod 4$ is good.
\end{lemma}

\begin{proof}
Strong induction on $n$. The case $n=1$ is done. If $n>1$ is odd, write
$n=pk$ with $p$ the smallest prime factor of $n$. Then $p$ is odd and $k<n$
is odd, so $n$ is good. If $n$ is even, then $4\mid n$. Write $n=4k$. If
$k\not\equiv 2\pmod 4$, use $k$. Otherwise $n=8m$ with $m=k/2$ odd, and we use $m$.
\end{proof}

\section{Charts for the building blocks}\label{sec:blocks}

All charts below are given in the data files of \texttt{general-q/data/charts}
(Section~\ref{sec:verify}).

\paragraph{$\F_4$ and $\F_8$.} The chart on $\F_4^5$ is the chart of
arXiv:2606.15860, \S A.6.2. The chart on $\F_8^5$ was found by simulated
annealing. For each of the $324$ reduced words, an inverse of $M_w$ is
recorded, and the products $M_wN_w=I$ are checked.

\paragraph{$\Z/p\Z$ for odd primes $p\le 79$.} For $p=3$ we use the chart of
arXiv:2606.15860, \S A.6.1. For $p=43$ we use the reduction modulo $43$ of
integer chart~1 below. Every other odd prime $p\le 79$ has a chart found by
simulated annealing. Each comes with inverse certificates as above.

\paragraph{Primes $p>79$.} We use six \emph{integer charts}: matrices
$A_a,A_b\in\GL_5(\Z)$ of determinant $\pm1$ with integer inverses, integer
vectors $b_a,b_b$, and $A_{t^{-1}}=A_t^{-1}$, $b_{t^{-1}}=A_{t^{-1}}b_t$, so
that~\eqref{eq:A1} holds over $\Z$. For every reduced word $w$, an integer
$d_w$ and an integer matrix $\mathrm{adj}_w$ are recorded with
$M_w\,\mathrm{adj}_w=d_wI$; here $d_w=\det M_w\ne 0$ and $\mathrm{adj}_w$ is the
adjugate. If a prime $p$ divides none of the $d_w$ of one chart, the
reduction of that chart modulo $p$ is universal by
Lemma~\ref{lem:transport}: $d_w$ is a unit of $\Z/p\Z$, and
$d_w^{-1}\mathrm{adj}_w$ is a right inverse of $M_w$. Let $P_i=\prod_w |d_w|$ for
chart $i$ and
\[
  \begin{aligned}
  g=\gcd(P_1,\dots,P_6)={}&2^{338}\,3^{146}\,5^{74}\,7^{34}\,11^{24}\,13^{22}\,17^{10}\,19^{12}\,23^{10}\\
  &\cdot 29^2\,31^2\,37^4\,41^2\,47^2\,53^2\,59^4\,61^2\,67^2\,71^2\,73^2\,79^4 .
  \end{aligned}
\]
A prime $p>79$ does not divide $g$, so it does not divide some $P_i$, and
chart~$i$ works modulo $p$. (The prime $43$ does not divide $g$ either.)

\begin{corollary}\label{cor:blocks}
$4$, $8$ and every odd prime are good.
\end{corollary}

\section{Proof of Theorem~\ref{thm:main}}\label{sec:proof}

By Corollary~\ref{cor:blocks} and Lemma~\ref{lem:reduce}, every $n\ge 1$ with
$n\not\equiv 2\pmod 4$ is good. Let $R$ be a ring with $n$ elements carrying
a universal chart, and let $C$ be the $144$-vertex controller, which is
exact-NB-$5$. By Lemma~\ref{lem:lift}, the lift graph on $C\times R^5$ has
$144\,n^5$ vertices and maximum degree at most $4n$, and every two vertices
are joined by a walk of length at most $5$. \qed

Explicitly, $R$ is a product of copies of $\F_4$ and $\F_8$ (at most one
$\F_8$) and of $\Z/p\Z$ for odd primes $p$ dividing $n$. For $n=q$ a prime
power one may instead take the field $\F_q$; this is
Theorem~\ref{thm:field}.

\section{The field version: proof of Theorem~\ref{thm:field}}\label{sec:field}

By Lemma~\ref{lem:transport} applied to the inclusion of a subfield, a
universal chart on $K^5$ is universal on $L^5$ for every field extension
$L\supseteq K$.

\begin{itemize}
\item $q=p^m$ with $p$ odd: $\F_q\supseteq\F_p$, and $\F_p$ carries a universal chart
by Section~\ref{sec:blocks}.
\item $q=2^m$ with $m$ even: $\F_q\supseteq\F_4$.
\item $q=2^m$ with $m$ odd and a prime factor $r\in\{3,5,7,11,13\}$:
$\F_q\supseteq\F_{2^r}$. Universal charts on $\F_8^5$, $\F_{32}^5$ and
$\F_{128}^5$ were found by simulated annealing, and on $\F_{2048}^5$ and
$\F_{8192}^5$ by random sampling. Each is checked on all $324$ reduced words.
\item $q=2^m$ with every prime factor of $m$ at least $17$: then
$q\ge 2^{17}$, and we use the Schwartz--Zippel lemma as follows.
\end{itemize}

Treat the $60$ entries of $A_a$, $A_b$, $b_a$, $b_b$ as indeterminates and
set $A_{t^{-1}}=\mathrm{adj}(A_t)/\det A_t$ and
$b_{t^{-1}}=\mathrm{adj}(A_t)b_t/\det A_t$ for $t\in\{a,b\}$. Multiply column $j$
of $M_w$ by $\det A_a$ once for each occurrence of $a^{-1}$ among
$t_{j+1},\dots,t_5$, and once more if $t_j=a^{-1}$; do the same for $b$. The
result has polynomial entries. Each factor $A_t$ contributes degree $1$ and
each $\mathrm{adj}(A_t)$ degree $4$, and the last factor $b_t$ or
$\mathrm{adj}(A_t)b_t$ has degree at most $5$. So column $j$ has degree at most
$4(5-j)+5$. The determinant $D_w$ of the modified matrix therefore has degree
at most $\sum_{j=1}^5\bigl(4(5-j)+5\bigr)=65$, and it equals $\det M_w$ times
powers of $\det A_a$ and $\det A_b$. The polynomial
\[
  P=\det A_a\cdot\det A_b\cdot\prod_{w}D_w
\]
has degree at most $10+324\cdot 65=21{,}070$ and integer coefficients. It does
not vanish at the $\F_4$ chart, so its reduction modulo $2$ is a nonzero
polynomial over $\F_2$. For $q=2^m>21{,}070$, the Schwartz--Zippel lemma gives
a point of $\F_q^{60}$ where $P\ne 0$. There $A_a$ and $A_b$ are invertible and
every $\det M_w\ne 0$, which is a universal chart on $\F_q^5$.

The graph statement follows from Lemma~\ref{lem:lift} and the Moore bound
remark after Corollary~\ref{cor:pp}. The graph is $4q$-regular: the
controller is simple, so distinct labels lead to distinct controller
vertices, and $b_t\ne 0$ because $b_t$ is the last column of the nonsingular
matrix $M_{ttttt}$. \qed

\section{The case $\F_2$: proof of Proposition~\ref{prop:two}}\label{sec:two}

Suppose $(A_t,b_t)$ is a universal chart on $\F_2^5$. Since $M_{aaaaa}$ is
nonsingular, $b_a$ is a cyclic vector of $A_a$, so $A_a$ is cyclic.
Simultaneous conjugation by $Q\in\GL_5(\F_2)$, that is
$(A_t,b_t)\mapsto(QA_tQ^{-1},Qb_t)$, preserves universality because
$M_w\mapsto QM_w$. So we may assume that $A_a$ is the companion matrix of its
characteristic polynomial $f$, of degree $5$ with $f(0)=1$ ($16$
possibilities). The centralizer of a cyclic matrix acts simply transitively on
its cyclic vectors, so we may further assume that $b_a=e_0$. An exhaustive
search over all $9{,}999{,}360$ matrices $A_b\in\GL_5(\F_2)$ and all $31$
nonzero $b_b$, for each $f$, finds no universal chart
(\texttt{scripts/f2exhaust.c}). Its evaluation of the $324$ words was
cross-checked against an independent implementation on random charts.

A finite commutative ring is a product of finite local rings, each of prime
power order. If $|R|\equiv 2\pmod 4$, exactly one factor has even order, and
that order is $2$, so the factor is $\F_2$. The projection $R\to\F_2$ would map
a universal chart on $R^5$ to one on $\F_2^5$ (Lemma~\ref{lem:transport}). \qed

\section{Data, verification and the Lean formalization}\label{sec:verify}\label{sec:lean}

The folder \texttt{general-q}, inside \texttt{results/diameter5-paper} of the
repository \url{https://github.com/rmizuno42/ddp} next to the materials of
arXiv:2606.15860, contains:

\begin{itemize}
\item \texttt{data/controller.json}: the arc maps $\sigma_t$ of the $144$-vertex
controller, exported from the paper's supplementary code, and for every
ordered pair of vertices a reduced word of length $5$ joining them.
\item \texttt{data/charts/}: the charts over $\F_4$, $\F_8$, $\Z/p\Z$
($3\le p\le 79$), the six integer charts, and the extension-field charts of
Section~\ref{sec:field}.
\item \texttt{scripts/}: independent Python checks of the charts, the
certificate generator, the search programs, and the $\F_2$ search.
\item \texttt{lean/}: the Lean~4 formalization of Theorem~\ref{thm:main}.
\end{itemize}

\paragraph{Lean.} The formalization (Lean $4.32.2$, Mathlib) defines
controllers, route charts over arbitrary commutative rings and lift graphs,
and proves Lemmas~\ref{lem:lift}--\ref{lem:reduce} and the argument for primes
$p>79$ in general. The finite facts are proved by \texttt{decide +kernel},
that is, evaluated by the Lean kernel itself:
\begin{itemize}
\item the controller identities $\sigma_{t^{-1}}\sigma_t=\mathrm{id}$, and the
certificate words for all $144^2$ ordered pairs;
\item for each chart, relations~\eqref{eq:A1} and, for every reduced word, the
certificate identity $M_wN_w=I$ (or $M_w\,\mathrm{adj}_w=d_wI$);
\item the factorization of $g$.
\end{itemize}
To keep the kernel computation small, the checks run on lists of
natural numbers or integers, and lemmas proved once show that this list
arithmetic agrees with the corresponding matrix operations over the ring.
$\F_4$ and $\F_8$ are defined by their addition and multiplication tables, and
their ring axioms are checked by \texttt{decide}. The main theorem is
\texttt{DDP.main} (and \texttt{DDP.main\_prime\_power} for prime powers). It has
no hypotheses beyond its statement, and \texttt{native\_decide} is not used.
A GitHub Actions workflow builds the whole project and checks with
\texttt{\#print axioms} that both theorems depend only on Lean's standard
axioms \texttt{propext}, \texttt{Classical.choice} and \texttt{Quot.sound}.
Theorem~\ref{thm:field} and
Proposition~\ref{prop:two} are not formalized.

\paragraph{The case $n=5$.} The lift over $\Z/5\Z$ (equal to $\F_5$) has $450{,}000$
vertices and $4{,}500{,}000$ edges. Its edge list was checked with the
repository's exact bit-parallel breadth-first search from every vertex:
it is $20$-regular, connected, and of diameter $5$. The same construction
program reproduces the published $(12,5)$ and $(16,5)$ edge lists of
arXiv:2606.15860 exactly.

\paragraph{Comparison.} The table of F.~Comellas covers $\Delta\le 16$. In
the CombinatoricsWiki degree--diameter table (last modified 18 February 2022)
the entry for $(20,5)$ is $281{,}820$. We have not compared with later results.
For $\Delta=4n\ge 28$ we know of no table, and we have not reviewed the
literature on large-degree constructions of diameter $5$; the family gives
$N(\Delta,5)\ge (144/1024)\,\Delta^5\approx 0.14\,\Delta^5$ for these $\Delta$.

\section*{Acknowledgements}
The controller and the charts over $\F_3$ and $\F_4$ come from
arXiv:2606.15860, where they were found through interaction with ChatGPT
(OpenAI). The extension in this note, the chart searches, the certificates and
the Lean formalization were produced with Claude Code (Anthropic) under the
author's direction.

\begin{thebibliography}{9}
\bibitem{M} R.~Mizuno, \emph{New lower bounds for the degree/diameter problem
via interaction with a browser-accessible LLM}, arXiv:2606.15860 (2026).
\bibitem{C} F.~Comellas, \emph{The (Degree, Diameter) Problem for Graphs},
\url{https://web.mat.upc.edu/francesc.comellas/delta-d/table_degree_diameter.html}.
\bibitem{CW} Combinatorics Wiki, \emph{The Degree Diameter Problem for General
Graphs}, \url{http://combinatoricswiki.org/wiki/The_Degree_Diameter_Problem_for_General_Graphs}.
\end{thebibliography}

\end{document}
